% =========================
% Explanation: Empirical Methodology and Study Type (informal, readable)
% =========================
\subsection*{Explanation of Empirical Methodology and Study Type (informal)}
This section explains what kind of empirical study we conduct, what it is designed to answer, and what claims it can (and cannot) support.

\begin{itemize}
    \item \textbf{Study type.}
    We treat the evaluation as a \emph{benchmarking / quantitative simulation} study of stochastic black-box optimization.
    Concretely, we compare multiple HPO/NAS algorithms by running each method on the same set of benchmark problem instances
    (datasets, splits, training recipes, and search spaces) under identical budgets and evaluation pipelines.
    Because objective evaluations are typically expensive and noisy (e.g., training a model), the study focuses on \emph{sample efficiency}
    measured as a function of the number of evaluations.

    \item \textbf{Core evaluation question.}
    The central question is whether the proposed method produces better outcomes \emph{under the same evaluation budget}
    (fixed-budget comparisons) and/or reaches a target performance level in fewer evaluations (fixed-target comparisons)
    than strong baselines.
    In the multiobjective case, ``better outcomes'' means recovering a higher-quality approximation to the Pareto front
    (e.g., higher hypervolume or lower \(\epsilon\)-indicator) under the same budget.

    \item \textbf{Hypotheses (pre-specified).}
    To avoid post-hoc selection of favorable outcomes, we pre-specify hypotheses:
    \begin{itemize}
        \item \textbf{Primary (efficiency):} Under identical budgets and evaluation pipelines, the proposed method achieves
        lower fixed-budget regret (single objective) and/or improved Pareto quality (multiobjective) compared to baselines.
        \item \textbf{Secondary (time-to-target):} The proposed method reaches near-oracle targets in fewer evaluations than baselines.
        \item \textbf{Mechanistic (attribution):} Any gains are attributable to the novel components of the method,
        verified via ablations (e.g., removing each new component and measuring the impact).
    \end{itemize}

    \item \textbf{Scope of claims.}
    We bound claims to what is actually tested:
    the specified search space \(\Lambda\), evaluation definition (what constitutes one evaluation),
    datasets/splits, training recipe, and the examined budget tiers.
    We do \emph{not} claim generalization to substantially different model families, datasets, training procedures,
    or evaluation models (e.g., multi-fidelity vs.\ full-fidelity) unless those settings are included in the benchmark suite.

    \item \textbf{Out-of-scope factors and interpretation.}
    Because this is a benchmarking study focused on evaluation counts, we define ``runtime'' primarily as the number of completed evaluations
    (trained models), not wall-clock time; wall-clock comparisons are only meaningful under tightly matched systems and parallelism.
    Likewise, we do not claim that the best observed configuration is a global optimum; instead, we evaluate relative to an empirical oracle reference
    constructed from the union of observed results at the largest tested budget.

    \item \textbf{Reproducibility framing.}
    We distinguish \emph{exact reproducibility} (possible for deterministic/tabular benchmarks or fully deterministic training stacks)
    from \emph{broad reproducibility} (typical for GPU training where nondeterminism remains).
    In both regimes, we specify the artifacts required to reproduce the reported figures and statistical conclusions:
    code, configurations, seeds, environment versions, and raw trial logs.
\end{itemize}
